% Звено повышенной частоты: структурная (блок-) схема преобразователя
% со звеном ВЧ — концепция А.В. Кобзева.
\documentclass[border=4pt]{standalone}
\usepackage{amsmath}
\usepackage{fontspec}
\setmainfont{Calibri}
\usepackage{tikz}
\usetikzlibrary{arrows.meta, positioning, calc}
\definecolor{ink}{RGB}{31,42,51}
\definecolor{grp}{RGB}{46,110,140}
\definecolor{light}{RGB}{233,247,245}

\begin{document}
\begin{tikzpicture}[
  line width=1.0pt, color=ink, font=\small,
  block/.style={draw=ink, fill=light, rounded corners=2pt,
                minimum height=1.05cm, inner sep=5pt, align=center},
  >={Stealth}]
  % --- верхний ряд
  \node[block] (set) {Сеть\\ $u_1$, $f_1$};
  \node[block, right=0.85 of set] (vyp) {Выпрямитель\\ + фильтр};
  \node[block, right=0.85 of vyp, fill=white, thick] (inv) {Инвертор ВЧ\\ (ШИМ)};
  \node[block, right=0.85 of inv, fill=white, thick] (tr) {ВЧ-трансфор-\\матор};
  % --- нижний ряд (змейкой)
  \node[block, below=1.0 of tr] (dem) {Демодулятор\\ (ВЧ-выпрями-\\тель)};
  \node[block, left=0.85 of dem] (fil) {Сглаживающий\\ фильтр};
  \node[block, left=0.85 of fil] (nag) {Нагрузка};
  % --- стрелки
  \draw[->] (set) -- (vyp);
  \draw[->] (vyp) -- (inv);
  \draw[->] (inv) -- (tr);
  \draw[->] (tr) -- (dem);
  \draw[->] (dem) -- (fil);
  \draw[->] (fil) -- (nag);
  % --- частота звена
  \node[text=grp, font=\footnotesize, anchor=west] at (inv.north west)
       {\strut};
  \node[text=grp, align=center, font=\footnotesize]
       at ($(inv.north)!0.5!(tr.north)$) (lbl) {};
  \node[text=grp, font=\footnotesize, below=0.32 of tr]
       {$f_{\text{п}} = 20\ldots100$ кГц};
  % --- развязка
  \node[text=grp, font=\footnotesize, above=0.18 of tr]
       {гальваническая развязка};
\end{tikzpicture}
\end{document}
